\documentclass[11pt]{article}
\usepackage[margin=1in]{geometry}
\usepackage{booktabs}
\usepackage{graphicx}
\usepackage{hyperref}
\usepackage{amsmath}
\usepackage{longtable}

\title{An Evidence-Backed Evaluation of CRME}
\author{}
\date{}

\begin{document}
\maketitle

\begin{abstract}
This manuscript presents an evidence-backed evaluation of the Cognitive Research Memory Engine (CRME), a structured research-memory architecture designed to organize and preserve research evidence, sessions, decisions, relations, and artifacts.
The evaluation was conducted across 2 datasets using 5 evaluation metrics and was followed by component-level ablation analysis.
The comparative evaluation identified SRCB\_v1 as the best-performing dataset across the reported metrics.
The findings are interpreted as descriptive and comparative evidence. No inferential statistical test or confidence interval was included in the current evidence package.
\end{abstract}

\section{Introduction}
Modern research workflows generate heterogeneous information, including experimental results, project decisions, sessions, relationships, and digital artifacts.
Without structured memory mechanisms, important research context can become fragmented across sessions and tools.
CRME addresses this problem through a structured research-memory architecture intended to preserve and organize research evidence while maintaining traceability between raw evaluation outputs and higher-level research artifacts.

\section{Evaluation Design}
The evaluation used 2 datasets and 5 evaluation metrics.
The evaluation pipeline consisted of experiment execution, dataset comparison, component ablation analysis, and statistical and descriptive evidence aggregation.

\begin{itemize}
\item MCS
\item SRS
\item KGQ
\item RRS
\item CES
\end{itemize}

\section{Results}
The comparative evaluation identified SRCB\_v1 as the best-performing dataset across all reported metrics.
\begin{table}[htbp]
\centering
\caption{Baseline Performance of CRME}
\label{tab:crme-baseline-performance}
\begin{tabular}{lc}
\toprule
Metric & Score \\
\midrule
MCS & 47.5 \\
SRS & 100 \\
KGQ & 40 \\
RRS & 43.33 \\
CES & 57.71 \\
\bottomrule
\end{tabular}
\end{table}

\section{Ablation Study}
The ablation analysis evaluated 4 architectural components under 5 configurations.
The observed ranking supports comparative component contribution analysis. It should not be interpreted as a definitive causal decomposition of architectural importance.
\begin{table}[htbp]
\centering
\caption{Component Contribution Based on Ablation}
\label{tab:crme-ablation-ranking}
\begin{tabular}{lcc}
\toprule
Rank & Component & Aggregate Degradation \\
\midrule
1 & decisions & 52.08 \\
2 & relations & 50 \\
3 & sessions & 20.83 \\
4 & artifacts & 9.38 \\
\bottomrule
\end{tabular}
\end{table}

\section{Statistical Analysis}
The current evidence package contains descriptive baseline metrics and component-level degradation values.
No inferential statistical test, confidence interval, or hypothesis test is currently included in the evidence package.
Accordingly, the present findings should be interpreted as descriptive and comparative rather than as statistically significant evidence of superiority over a baseline system.

\section{Discussion}
The results indicate heterogeneous contribution across the evaluated CRME components.
The observed ranking should be understood as empirical comparative evidence rather than definitive causal estimates.

\section{Threats to Validity}
The current evidence package does not include inferential statistical testing or confidence interval estimation.
The ablation results provide comparative evidence of component contribution but do not establish definitive causal importance.
Generalization to other datasets and research workflows requires further validation.

\section{Reproducibility}
All reported results are derived from stored CRME evaluation artifacts.
The evaluation pipeline preserves the relationship between raw evidence, aggregated summaries, paper tables, figures, and claim-level evidence.
The Q1 consistency audit completed with overall status \texttt{PASS}.
Raw evaluation outputs are preserved and are not modified by the paper generation layer.

\section{Conclusion}
This evidence-backed evaluation establishes a reproducible descriptive baseline for CRME across the currently evaluated datasets and metrics.
The comparative and ablation results indicate heterogeneous observed contribution across the evaluated architectural components.
Stronger claims regarding statistical superiority or definitive causal importance require additional experimental and inferential validation.

\end{document}